%! Logarithmic Functions — Unit 2, Lesson 2.11 — Student Packet Cover Sheet
\documentclass[10pt]{article}
\usepackage{precalculus-article}
\usepackage{precalculus-boxes}
\usepackage{ltablex}
\keepXColumns

\begin{document}

% Full-bleed navy banner
\begin{tikzpicture}[remember picture, overlay]
  \fill[navy] (current page.north west)
    rectangle ([yshift=-0.9in]current page.north east);
\end{tikzpicture}
\vspace{-0.8in}

\begin{center}
  {\color{white}\textbf{\LARGE AP Precalculus}}\\[4pt]
  {\color{skymid}\large\textbf{Unit 2: Exponential and Logarithmic Functions}}\\[6pt]
  {\color{white}\textbf{\large Lesson 2.11 \quad Logarithmic Functions}}
\end{center}

\vspace{0.22in}
\namedateperiod
\vspace{0.18in}

\begin{learningtargetbox}
\small
\begin{enumerate}[leftmargin=1.5em, itemsep=5pt]
  \item I can identify the domain, range, vertical asymptote, and key characteristics of a
        logarithmic function in general form $f(x) = a\log_b x$.
        \hfill\textit{LO 2.11.A}
  \item I can describe the end behavior and monotonicity of a logarithmic function, and
        explain how they connect to the inverse exponential function.
  \item I can use the additive transformation identification test to determine whether a
        function is logarithmic.
\end{enumerate}
\end{learningtargetbox}

\vspace{0.14in}

\begin{tocbox}
\small
\renewcommand{\arraystretch}{1.5}
\begin{tabularx}{\linewidth}{@{} c l X r @{}}
\rowcolor{sky}
\textbf{\#} & \textbf{Component} & \textbf{Description} & \textbf{Score} \\
\midrule
1 & Warm-Up        & Spiral review: evaluating logarithms and domain restrictions & \blank{1.2cm} \\
2 & Guided Notes   & Domain, range, asymptote, end behavior, concavity, identification test & \blank{1.2cm} \\
3 & Group Activity & Analyzing characteristics of logarithmic functions by tier & \blank{1.2cm} \\
4 & Exit Ticket    & Domain/asymptote and end-behavior checks & \blank{1.2cm} \\
5 & Homework       & Independent practice with characteristics of logarithmic functions & \blank{1.2cm} \\
\midrule
  & \multicolumn{2}{r}{\textbf{Total}} & \blank{1.2cm} \\
\end{tabularx}
\end{tocbox}

\end{document}
